\documentclass[11pt,a4paper]{article}

\usepackage[spanish,es-noshorthands]{babel}
\usepackage[T1]{fontenc}
\usepackage[utf8]{inputenc}
\usepackage[a4paper,margin=2.4cm]{geometry}
\usepackage{amsmath,amssymb}
\usepackage{listings}
\usepackage{tikz}
\usetikzlibrary{arrows.meta,positioning}
\usepackage{xcolor}
\usepackage[colorlinks=true,linkcolor=blue,urlcolor=blue]{hyperref}

\lstset{
  basicstyle=\ttfamily\small,
  frame=single,
  breaklines=true,
  keepspaces=true,
  columns=fullflexible,
  language={},
  keywordstyle=\color{black}
}

\title{\textbf{Modelado en VDM++ de los Componentes del Subsistema de Rutas}}
\author{Proyecto: Sistema de Planificación de Viajes en Transporte Público}
\date{\today}

\begin{document}

\maketitle

\begin{abstract}
Este documento presenta el modelado inicial en VDM++ del subsistema de rutas del proyecto de planificación de viajes en transporte público. Se identifican los tipos \emph{quote} y \emph{token}, se modela un registro con una invariante, se analiza el uso de los tipos secuencia (\texttt{seq}) y conjunto (\texttt{set}), se propone un diagrama de clases inicial y se utiliza el tipo \texttt{map} para relacionar dichas clases, indicando la invariante asociada.
\end{abstract}

\section{Tipos Quote y Token}

Un tipo \textbf{quote} representa un conjunto finito de valores simbólicos (una enumeración). En el subsistema de rutas, el campo \texttt{sentido} del recorrido solo admite dos valores, por lo que se modela como un tipo quote:

\begin{lstlisting}
Sentido = <IDA> | <VUELTA>
\end{lstlisting}

Un tipo \textbf{token} representa valores opacos que no admiten ninguna operación, salvo la igualdad; se usa para identificadores únicos cuyo contenido interno no es relevante para la especificación:

\begin{lstlisting}
RutaId     = token
ParaderoId = token
\end{lstlisting}

\section{Registro modelado con invariante}

Un \textbf{registro} (\emph{record}) es un tipo compuesto por campos con nombre. Se modelan dos registros: \texttt{Punto}, que representa una coordenada geográfica, y \texttt{Recorrido}, que agrupa los datos de un trayecto.

\begin{lstlisting}
Punto :: lat : real
         lng : real;

Recorrido :: ruta         : RutaId
             sentido      : Sentido
             coordenadas  : seq of Punto
             color        : seq of char;
\end{lstlisting}

Se definen dos \textbf{invariantes}. La primera garantiza que las coordenadas estén dentro de un rango geográfico válido; la segunda, que todo recorrido tenga al menos dos puntos, coincidiendo con la regla de negocio del sistema.

\begin{lstlisting}
inv punto ==
  lat >= -90  and lat <= 90  and
  lng >= -180 and lng <= 180;

inv recorrido ==
  len coordenadas >= 2;
\end{lstlisting}

\section{Uso de Secuencia (Seq) y Conjunto (Set)}

El proyecto hace uso de \textbf{ambos} tipos: secuencia y conjunto.

\begin{itemize}
  \item \textbf{Secuencia (\texttt{seq})}: el campo \texttt{coordenadas : seq of Punto}. Modela la propiedad de \emph{orden} del trayecto: el orden de los puntos define el recorrido y es posible que un punto se repita (por ejemplo, en una parada). Por ello se elige secuencia y no conjunto.
  \item \textbf{Conjunto (\texttt{set})}: colecciones como \texttt{paraderos : set of ParaderoId} o \texttt{empresas : set of Empresa}. Modelan catálogos en los que el \emph{orden no importa} y \emph{no debe haber duplicados}, por lo que el conjunto es la representación natural.
\end{itemize}

\section{Diagrama de clases inicial}

El diagrama incluye las clases \texttt{Empresa}, \texttt{Ruta}, \texttt{Recorrido}, \texttt{Paradero} y \texttt{RecorridoParadero}, junto con sus asociaciones y cardinalidades.

\begin{center}
\begin{tikzpicture}[
  class/.style={draw, rectangle, rounded corners=2pt, align=left, inner sep=5pt, font=\small},
  assoc/.style={-{Stealth[length=2mm]}, thick},
  every node/.style={}
]

\node[class] (empresa) at (0,10.5) {
  \begin{tabular}{l}
    \textbf{Empresa}\\
    \hline
    nombre : seq of char\\
    descripcion : seq of char
  \end{tabular}
};

\node[class] (ruta) at (0,7.5) {
  \begin{tabular}{l}
    \textbf{Ruta}\\
    \hline
    nombre : seq of char\\
    codigo : seq of char
  \end{tabular}
};

\node[class] (recorrido) at (-2.5,4) {
  \begin{tabular}{l}
    \textbf{Recorrido}\\
    \hline
    sentido : Sentido\\
    coordenadas : seq of Punto
  \end{tabular}
};

\node[class] (paradero) at (4.5,4) {
  \begin{tabular}{l}
    \textbf{Paradero}\\
    \hline
    nombre : seq of char\\
    latitud : real\\
    longitud : real
  \end{tabular}
};

\node[class] (rpar) at (1,1) {
  \begin{tabular}{l}
    \textbf{RecorridoParadero}\\
    \hline
    orden : nat\\
    distancia : real
  \end{tabular}
};

\draw[assoc] (empresa.south) -- (ruta.north)
  node[pos=0.55,left] {1} node[pos=0.55,right] {$*$};

\draw[assoc] (ruta.south) -- (recorrido.north)
  node[pos=0.55,left] {1} node[pos=0.55,right] {1..2};

\draw[assoc] (recorrido.south) -- (rpar.north west)
  node[pos=0.5,left] {1} node[pos=0.5,right] {$*$};

\draw[assoc] (paradero.south) -- (rpar.north east)
  node[pos=0.5,left] {$*$} node[pos=0.5,right] {1};

\end{tikzpicture}
\end{center}

\section{Tipo MAP para relacionar las clases}

El tipo \texttt{map} relaciona un dominio con un rango. Para vincular cada ruta con sus recorridos se define:

\begin{lstlisting}
recorridosPorRuta : map Ruta to set of Recorrido
\end{lstlisting}

La \textbf{invariante asociada} establece dos propiedades: (1) cada ruta tiene a lo sumo dos recorridos (IDA y VUELTA), y (2) cada recorrido almacenado en el mapa hace referencia a la ruta que lo contiene.

\begin{lstlisting}
inv recorridosPorRuta ==
  (forall r in set dom recorridosPorRuta &
     card recorridosPorRuta(r) <= 2)
  and
  (forall r in set dom recorridosPorRuta &
     forall rec in set recorridosPorRuta(r) &
       rec.ruta = r);
\end{lstlisting}

\end{document}
